\documentclass[11pt]{article}

\usepackage[T1]{fontenc}
\usepackage[margin=1in]{geometry}
\usepackage{xcolor}
\usepackage{listings}
\usepackage{tabularx}
\usepackage{hyperref}

\definecolor{codegray}{RGB}{248,248,248}
\lstset{
  basicstyle=\small\ttfamily,
  backgroundcolor=\color{codegray},
  frame=single,
  breaklines=true,
  columns=fullflexible,
  keepspaces=true,
  showstringspaces=false
}

\title{ESBMC Initialization Display: Before and After the Fix}
\author{ESBMC}
\date{October 22, 2026}

\begin{document}
\maketitle

\section{Summary}

The listings below compare counterexample excerpts from the same programs and
commands before and after the repair. State headers and unrelated trace steps
are omitted. Both examples intentionally violate their assertions and report
\texttt{VERIFICATION FAILED} in both versions.

\begin{tabularx}{\textwidth}{|l|X|X|}
\hline
\textbf{Category} & \textbf{Before the fix} & \textbf{After the fix} \\
\hline
Nested structure & The nested member remains symbolic despite constrained leaf values. & Nested fields show concrete values: \texttt{i=7}, \texttt{d=1.25}. \\
\hline
Multidimensional array & A two-by-two array displays eight flattened elements. & Four elements appear as two rows of two elements. \\
\hline
\end{tabularx}

\section{Nested Structure}

Save as \texttt{nested.c}:

\begin{lstlisting}[language=C]
struct Inner
{
  int i;
  double d;
};

struct Outer
{
  struct Inner inner;
  int values[2];
};

int main(void)
{
  struct Outer obj;
  __ESBMC_init_object(&obj);
  __ESBMC_assume(obj.inner.i == 7 && obj.inner.d == 1.25 &&
                 obj.values[0] == 3 && obj.values[1] == 4);
  __ESBMC_assert(!(obj.inner.i == 7 && obj.inner.d == 1.25 &&
                   obj.values[0] == 3 && obj.values[1] == 4),
                 "audit failure");
  return 0;
}
\end{lstlisting}

\begin{lstlisting}[language=bash]
esbmc nested.c --z3
\end{lstlisting}

\subsection*{Before the fix}

\begin{lstlisting}
obj = { .inner=nondet_symbol(symex::nondet1).inner,
        .values={ 3, 4 } }
\end{lstlisting}

\subsection*{After the fix}

\begin{lstlisting}
obj = { .inner={ .i=7, .d=1.250000e+0 }, .values={ 3, 4 } }
\end{lstlisting}

\noindent\textbf{Visible change:} The symbolic expression
\texttt{nondet\_symbol(symex::nondet1).inner} is replaced by the nested value
\texttt{\{ .i=7, .d=1.250000e+0 \}}. The array member still displays
\texttt{\{ 3, 4 \}}. The floating-point notation represents the constrained
value \texttt{1.25}.

\section{Multidimensional Array}

Save as \texttt{multidimensional.c}:

\begin{lstlisting}[language=C]
int main(void)
{
  int obj[2][2];
  __ESBMC_init_object(&obj);
  __ESBMC_assume(obj[0][0] == 1 && obj[0][1] == 2 &&
                 obj[1][0] == 3 && obj[1][1] == 4);
  __ESBMC_assert(!(obj[0][0] == 1 && obj[0][1] == 2 &&
                   obj[1][0] == 3 && obj[1][1] == 4),
                 "audit failure");
  return 0;
}
\end{lstlisting}

\begin{lstlisting}[language=bash]
esbmc multidimensional.c --z3
\end{lstlisting}

\subsection*{Before the fix}

\begin{lstlisting}
obj = { 1, 2, 3, 4, 2, 2, 2, 2 }
\end{lstlisting}

\subsection*{After the fix}

\begin{lstlisting}
obj = { { 1, 2 }, { 3, 4 } }
\end{lstlisting}

\noindent\textbf{Visible change:} Eight flattened entries become four elements
grouped by their original dimensions. The first row is \texttt{\{ 1, 2 \}} and
the second row is \texttt{\{ 3, 4 \}}; the four extra entries disappear.

\section{Implementation and Solver Validation}

The repair uses the common SMT model-extraction layer and the symbol tuple
flattener. Structured values are read through the tuple AST interface, and
finite multidimensional arrays are rebuilt using their source dimensions and
bounds. The implementation is shared across solver backends.

\subsection{Solver Validation}

The two examples and their successful property counterparts were checked with
each backend below. Each table entry covers four runs.

\begin{tabularx}{\textwidth}{|l|X|X|X|}
\hline
\textbf{Backend} & \textbf{Default} & \textbf{Node flattener} & \textbf{Symbol flattener} \\
\hline
Z3 & 4/4 passed & 4/4 passed & 4/4 passed \\
\hline
Boolector & 4/4 passed & 4/4 passed & 4/4 passed \\
\hline
CVC5 & 4/4 passed & 4/4 passed & 4/4 passed \\
\hline
Yices & 4/4 passed & 4/4 passed & 2/4 passed \\
\hline
Bitwuzla & 4/4 passed & 4/4 passed & 4/4 passed \\
\hline
\end{tabularx}

All five backends pass in their default configuration. With Yices and the
symbol flattener, the nested-structure cases fail during SMT encoding with
\texttt{argument is not a bitvector}, before model extraction. That combination
is not counted as supported by this repair.

The SMT-LIB path also displayed the nested structure correctly with external
Z3, CVC5, Yices, and Bitwuzla processes. Native CVC4 and MathSAT backends were
not tested.

\end{document}
